\section{Interface stability and exact interactions}\label{sec:interface}

We retain the notation and hypotheses of Section~\ref{sec:reach}.  Thus
$n=7$, $m=6$, and the penalty $\Gamma$, including its constant value
$G_1$ near $1$, is fixed.  None of the estimates in this section assumes
an energy density bound.  The elementary local estimates below also
hold in every fixed ambient dimension.

Our objective is to turn ambient stability into estimates on the zero
surface and to identify the strength of the interaction between nearby
zero sheets.  Recall that $D=1-|\nabla U|^2$ is a diffuse scalar field,
whereas $a^2=|A|^2$ and $H=\tr A$ are, respectively, the squared norm
and signed trace of the second fundamental form of a regular zero sheet.
These quantities have different roles: the estimates below compare
them, but do not identify them or assume that $H$ vanishes.

The argument has three parts.  A finite normal average of $q^3D$
controls curvature and mean curvature, and supplies an additional
favorable term in surface stability.  A comparison on both sides of a
sheet then relates $D$ to curvature and the exponential of its penalized
reach.  Finally, on logarithmic windows with uniformly small defect,
we determine the sharp constants in the interaction and phase-tail
estimates.  These are the surface inputs for the growth and radial
arguments in Sections~\ref{sec:growth} and~\ref{sec:radial}.
The passage from stable diffuse solutions to interacting transition
layers is central to the local theory of \citet{WangWei2019}.  Here we
prove the surface test and exact interaction estimates needed in the
quartic normalization.

\subsection{Normal averages, curvature, and mean curvature}

Recall the normal average defined in
\eqref{eq:diffuse-normal-average}: at a regular zero and for $P>0$,
\[
 I_P(z)=\frac1\sigma\int_{-P}^{P}
 q(z+s\nu(z))^3D(z+s\nu(z))\,\dd s,
 \qquad \sigma=\frac{4\sqrt2}{3}.
\]
Once $P$ has been fixed we write $I=I_P$.
The normalization uses $\sigma=\int_\R G^2\,\dd s$, rather than the
energy of one transition, $s_0=\sigma/2$.  In particular, $I_P$ is a
normal average based at $z$, not the pointwise defect $D(z)$.

\begin{lemma}[Sharp means]\label{lem:interface-sharp-means}
For each $\eps>0$ there exist $P_\eps<\infty$ and, for each fixed
$P\ge P_\eps$, a number $\delta_{\eps,P}>0$ such that every zero with
$0<D(z)<\delta_{\eps,P}$ is regular and satisfies
\begin{equation}\label{eq:L1}
 cD(z)\le I_P(z)\le C_P D(z),\qquad
 I_P(z)\ge(1-\eps)|A(z)|^2,\qquad
 |H(z)|\le\left(\frac1{2\sqrt2}+\eps\right)I_P(z).
\end{equation}
The constants and thresholds are independent of the solution and center.
The constant $c>0$ can be chosen independently of $P\ge1$.
\end{lemma}

\begin{proof}
It suffices to treat $0<\eps<1$, since decreasing $\eps$ strengthens
the conclusions.  The two-sided comparison with $D$ follows from the local Harnack bound in
\eqref{eq:diffuse-local}, the positive lower bound for $q$ in a fixed
small neighborhood of a zero, and its upper bound on a fixed segment.
The mean-curvature estimate is the finite-window conclusion of
Lemma~\ref{lem:diffuse-mean}.  To prove the additional curvature
estimate, and to identify the common limiting quantities, consider the
normalized limits of Lemma~\ref{lem:diffuse-compactness},
with $z=0$, $\nu(z)=e_n$, and normalization by $D(0)$ for $D$ and by
$\sqrt{D(0)}$ for $T$.  Write these limits as $\mathcal D,\mathcal T$.
Then $\mathcal D(0)=1$, $\mathcal T$ is a horizontal harmonic tensor
independent of $s=x_n$, and
\[
 F(x',s)=G(s)\mathcal D(x',s),\qquad
 (-\Delta+2-G)F=2G(s)|\mathcal T(x')|^2.
\]
We use the nonnegative source part of
Lemma~\ref{lem:diffuse-positive} to compare this mass with curvature.
The function $|\mathcal T|^2$ is horizontally subharmonic.  In a product cylinder
$B_R^m\times(-R,R)$ the positive Dirichlet Green kernel, evaluated at
$(0,s)$, is radial in the horizontal source variable.  Horizontal
spherical averaging therefore shows that its integral against
$2G|\mathcal T|^2$ is at least its integral against
$2G|\mathcal T(0)|^2$.  The zero-boundary source solution is bounded
above by $F$.  As $R\to\infty$, the latter constant-horizontal source
solutions increase to $|\mathcal T(0)|^2 f(s)$, where
\[
 (-\partial_s^2+2-G)f=2G,\qquad f\ge0,
\]
with $f$ bounded.  To justify this identification, comparison with
the constant $2$ bounds the exhaustions; every bounded entire
solution with this source is unique, since the potential is at least
$1$, and horizontal translations then force independence of $x'$.
The uniqueness follows by comparison on balls with the growing
exponential barriers for $-\Delta+1$.
Since $(-\partial_s^2+2-G)G=2G^2$, integration against $G$ gives
\[
 \frac1\sigma\int_\R G(s)^2F(0,s)\,\dd s
       \ge |\mathcal T(0)|^2.
\]
All integrations are legitimate by boundedness of $f$ and exponential
decay of $G$; one may first insert compact cutoffs.

These limiting inequalities imply their stated finite-window versions.
Here is the uniformity argument for the curvature assertion, where an
additive limiting error would be insufficient.  If no fixed interval
and threshold worked, choose $P_j\to\infty$, $D(z_j)\to0$, and
$I_{P_j}<(1-\eps)a(z_j)^2$.  The lower local mass bound gives
$a(z_j)^2/D(z_j)\ge c/(1-\eps)>0$.  Subsequence convergence on every
fixed interval would contradict the preceding limiting inequality.
Here $a^2/D$ converges to $|\mathcal T(0)|^2$.
Choose the interval large enough also for Lemma~\ref{lem:diffuse-mean}
with coefficient $1/(2\sqrt2)+\eps$, and then take the smaller of the
two defect thresholds.  Its nonnegative integrand preserves that
mean-curvature bound on every larger fixed interval.  This proves
all assertions with the stated order of choices.
\end{proof}

The next estimate concerns tangential variation of $H$.  Its small
coefficient is obtained by averaging on a sufficiently large, but fixed,
surface disk.  It will allow the mean-curvature derivative in the
curvature equation to be absorbed in Section~\ref{sec:growth}.

\begin{lemma}[Averaged mean derivative]\label{lem:interface-mean-gradient}
For every $\eps>0$ there are $S\ge10$ and $\delta>0$ such that, if
$z\in\Sigma$ and $D(z)<\delta$, all zeros in $B_{4S}(z)$ belong to
one regular graph converging smoothly to the tangent plane as the
threshold decreases, and
\begin{equation}\label{eq:L2}
 |\nabla_\Sigma H(z)|^2\le
 \frac{\eps}{S^m}\int_{B_S(z)\cap\Sigma_{\rm reg}}D^2\,\dd S.
\end{equation}
\end{lemma}

\begin{proof}
Smooth single-profile convergence on every fixed ball follows from
\eqref{eq:diffuse-local}.  In the normalized limit used above,
$\nabla_\Sigma H/D(z)$ converges to $\nabla_{x'}F_s(0)/2$.
The even source part of $F$ contributes zero to this derivative.
For a positive homogeneous extremal
$F=e^{b\cdot x'}f(s/\sqrt2)$, the explicit profiles in
Lemma~\ref{lem:diffuse-positive} give
\[
 |\nabla_{x'}F_s(0)|\le C|b|F(0),\qquad |b|\le\sqrt{3/2}.
\]
The average of $e^{b\cdot x'}$ on $B_{S/2}^m$ is at least $1$,
and tends to infinity uniformly when $|b|$ is bounded below by any
positive number.  First restricting to $|b|<\eta$, then taking $S$
large on its complement, proves, with any prescribed $\zeta>0$,
\[
 \frac12|\nabla_{x'}F_s(0)|
 \le \zeta\fint_{B_{S/2}^m}F(x',0)\,\dd x'.
\]
The same inequality holds for positive mixtures by the triangle
inequality; the nonnegative even source can only enlarge its right
side.  Choose $\zeta$ in terms of $\eps$ and apply Cauchy--Schwarz.
Since $F(x',0)=\mathcal D(x',0)$, this proves the limiting version
of \eqref{eq:L2}.  Smooth normalized convergence and the uniformly
positive local mass of $\mathcal D^2$ allow a fixed strict relaxation.
The graph area elements converge to $\dd x'$, and the disk
$B_{S/2}^m$ lies in the surface ball used on the right.  A contradiction
sequence now gives a uniform positive threshold $\delta$.
\end{proof}

\subsection{From ambient tests to surface stability}

The preceding estimates concern one zero and a fixed normal interval.
We now use them simultaneously along a surface test.  Normal segments
from different bases may overlap, so the ambient test must account for
all their inverse branches.  Taking the square root of the sum of their
squares preserves the potential term exactly and produces a useful
nonnegative difference of gradient energies.

\Needspace{6\baselineskip}
\begin{lemma}[Pushing a lifted test to space]\label{lem:interface-immersed-test}
Let $\Omega\subset\Sigma_{\rm reg}$ be open and
$\Phi(z,s)=z+s\nu(z)$.  Fix $0<\alpha<1/4$ and an even smooth
function $\beta(z,s)=\beta_0(sa(z))$, equal to $1$ for
$|s|a\le\alpha/2$ and supported in $|s|a<\alpha$.
For a compactly supported Lipschitz amplitude on the lifted tube,
write $\xi_j$ for its local inverse branches and put
$\xi=(\sum_j\xi_j^2)^{1/2}$.  Then $\xi$ is an admissible physical
test.  Its quadratic form is the lifted form minus
$\int\mathcal V\,\dd x$, where the nonnegative variance density is
\[
 \mathcal V=\min_{b\in\R^n}\sum_j|\nabla\xi_j-\xi_j b|^2.
\]
For the lifted amplitude $G(s)\beta(z,s)\psi(z)$, the lifted form is
\begin{equation}\label{eq:L4}
 \int_{\Omega\times\R}\left\{
 G^2J|\nabla(\beta\psi)|^2+
 \beta^2\psi^2\left[-J_sGG'
       +3\bigl(v(\Phi)^2-g(s)^2\bigr)G^2J\right]\right\}
       \,\dd S\,\dd s,
 \qquad J=\det(\Id+sA).
\end{equation}
Here the gradient and volume use the pullback Euclidean metric.
Formula \eqref{eq:L4} also holds when $\psi$ includes a bounded
Lipschitz factor depending on $x=\Phi(z,s)$, initially with compact
support and a finite normal cap.
\end{lemma}

\begin{proof}
On $|s|a<\alpha$, $\Phi$ is a local diffeomorphism and $J>0$.
A fiber restricted to the compact lifted support is discrete and
compact, hence finite.  A sufficiently small physical neighborhood
has only these nonzero inverse branches: otherwise a sequence of
additional branches has a limit in the same compact support and
contradicts local invertibility.  Consequently the squared sum is
locally a finite sum.  Its square root is Lipschitz, including at
its zero set, and almost everywhere
\[
 \sum_j|\nabla\xi_j|^2-|\nabla\xi|^2
 =\sum_j|\nabla\xi_j|^2-
   \frac{|\sum_j\xi_j\nabla\xi_j|^2}{\sum_j\xi_j^2}
 =\mathcal V\ge0.
\]
At points where all terms vanish, their gradients vanish almost
everywhere, so the same equality holds with value zero.  The potential
terms add exactly, irrespective of their sign.  The area formula for
$\Phi$ thus proves the assertion without an injective global tube.

For any lifted amplitude $G f$, expand its normal derivative and
integrate the cross term $2GG'ff_sJ$ by parts.  Using
$G''=(3g^2-1)G$ gives precisely \eqref{eq:L4}, with $f=\beta\psi$.
There are no artificial chart boundaries in this computation: the
amplitude is globally defined on the lifted surface, while local
inverse charts are used only to justify its pushdown.  Finite even
normal caps may be removed whenever the displayed terms are dominated
by a polynomial in $|s|$ times $G^2$; this is the case in the applications
below.  At $a=0$ take $\beta=1$ and use the same capped limit.
\end{proof}

Physical stability and the variance identity show that the lifted
quadratic form is nonnegative.  We may discard the subtracted variance
for the next estimate; Section~\ref{sec:growth} will retain it in a
more general gluing construction.  The next proposition combines two
negative terms in the lifted quadratic form.  The normal Jacobian
supplies $a^2/2$, while the profile deficit
supplies $I_P/2$.  Keeping both terms will allow separate control of
curvature and of interaction in the later estimates.

\begin{proposition}[Strengthened surface stability]
\label{prop:interface-strengthened-stability}
For each $\eps>0$, choose $P$ sufficiently large as in
Lemma~\ref{lem:interface-sharp-means}.  There exists $\delta>0$
such that every compactly supported Lipschitz function on
$\{z\in\Sigma: D(z)<\delta\}$ satisfies
\begin{equation}\label{eq:L3}
 (1-\eps)\int_\Sigma\frac{a^2+I_P}{2}\psi^2\,\dd S
 \le\int_\Sigma|\nabla_\Sigma\psi|^2\,\dd S.
\end{equation}
In the lifted computation the gradient coefficient can be made
arbitrarily close to $\sigma$, and the negative coefficient
arbitrarily close to $\sigma(a^2+I_P)/2$, with the same choices.
\end{proposition}

\begin{proof}
Use Lemma~\ref{lem:interface-immersed-test}.  On the support of $\beta$
the pullback tangential metric differs from the surface metric by
$O(\alpha)$.  Uniform derivative estimates give
$|\nabla\beta|\le C_\alpha(1+|s|)$; derivatives of $a=|A|$ are
understood weakly at its zero set.  All terms containing a derivative
of $\beta$ are bounded, after a Young inequality with any fixed small
parameter, by $C_\alpha e^{-c_\alpha/a}\psi^2$.  Thus the gradient
part is at most
\[
 (1+C\alpha+\eta)\sigma|\nabla_\Sigma\psi|^2
       +C_{\alpha,\eta}e^{-c_\alpha/a}\psi^2.
\]
The exponential remainder is $o(D)$ uniformly as $D\to0$, since
$a^2\le CD$.

The elementary-symmetric expansion of $J$ and evenness give
\[
 -\int\beta^2J_sGG'\,\dd s
   =\frac\sigma2(H^2-a^2)+o(D).
\]
Indeed the constant term $H$ in $J_s$ integrates to zero,
the linear term is $(H^2-a^2)s$, and
$-\int sGG'=\sigma/2$.  Terms of order at least two in $sA$
are either odd under the integral or bounded by $Ca^3=o(D)$.
The cap errors have the preceding exponential bound, and $H^2=o(D)$.

The potential contribution is nonpositive everywhere: $U(z)=0$
and $|\nabla U|\le1$ imply $|U(z+s\nu)|\le|s|$ and hence
$v(z+s\nu)^2\le g(s)^2$.  We quantify its gain.
On the two normal rays define
$U_\pm(r)=\pm U(z\pm r\nu)$ for $r\ge0$.  Then
\[
 r-U_\pm(r)\ge\frac12\int_0^rD(z\pm t\nu)\,\dd t.
\]
Indeed $U_\pm'(r)=p(z\pm r\nu)\cdot\nu
\le\sqrt{1-D(z\pm r\nu)}\le1-D(z\pm r\nu)/2$.
This writes the same deficit inequality on both sides without changing
its sign convention.
Moreover, on every fixed segment,
\begin{equation}\label{eq:interface-normal-accuracy}
 |p(z+s\nu)-\nu|+|U(z+s\nu)-s|\le C_P D(z),\qquad |s|\le P.
\end{equation}
To prove it, initially $|p(z)-\nu|\le CD(z)$, and along the segment
\[
 |T\nu|\le|Tp|+|T||p-\nu|
       \le C_P D(z)+C_P\sqrt{D(z)}|p-\nu|.
\]
A Gronwall estimate gives \eqref{eq:interface-normal-accuracy}.
Here $Tp=-\nabla D/2$ and \eqref{eq:diffuse-local} were used.

For a fixed segment $[-Q,Q]$, \eqref{eq:interface-normal-accuracy}
gives $p\cdot\nu>0$ once the threshold is small enough.  Thus $U(s)$
has the sign of $s$ on both half-segments.  After division by $D(z)$
and passage to a normalized limit, the
potential gain $3(g^2-v^2)G^2$ is at least
\[
 -\frac12(G^3)'(s)\int_0^s\mathcal D(0,t)\,\dd t.
\]
This expression is nonnegative on both half-lines.  Integrating
and using Tonelli, its total mass is
$\tfrac12\int_\R G^3\mathcal D\,\dd s$.
To make the finite-window order explicit, first fix $P$ and put
$h_P(s)=\int_0^s\mathbf1_{\{|t|\le P\}}\mathcal D(0,t)\,\dd t$.
For $Q\ge P$, integration by parts gives exactly
\[
 -\frac12\int_{-Q}^{Q}(G^3)'h_P\,\dd s
 =\frac12\int_{-P}^{P}G^3\mathcal D\,\dd s
       -\frac{G(Q)^3}{2}\int_{-P}^{P}\mathcal D\,\dd s.
\]
Local Harnack bounds the last integral by $C_P$, uniformly for
normalized limits.  Choose $Q$ large after $P$ so the loss is as
small as prescribed, then decrease the defect threshold to use
compactness on $[-Q,Q]$.  This threshold also ensures $\beta=1$
on that segment.  Since $I_P/D\ge c$, for any $\eta>0$ this gives
\[
 \int 3(g^2-v(\Phi)^2)G^2\beta^2J\,\dd s
       \ge(1-\eta)\frac\sigma2 I_P
\]
when the defect threshold is sufficiently small.  Fix $\eta$, then
$\alpha$ and the threshold, in that order, to absorb the $o(D)$ terms
using $I_P\ge cD$.  Physical stability and the upper bound by the
lifted form now prove \eqref{eq:L3}.
All normal intervals used here are fixed before the final defect
threshold: choose $P$, then the larger extraction interval $Q$, and
then decrease the threshold.  No later surface collar or surface-test
amplitude enters those choices.
\end{proof}

\subsection{Comparison at the exact exponential rate}

For $0<\delta\le1/4$ set
\[
 X=\{z\in\Sigma:D(z)\ge\delta\},\qquad
 \bar a(z)=\begin{cases}1,&z\in X,\\ a(z),&z\notin X,\end{cases}
 \qquad \kappa=\frac{\sqrt2}{G_1},\qquad E(z)=e^{-\kappa d(z)}.
\]
The zero set is closed, and $X$ includes every singular zero.
Here $E$ records the interaction scale associated with reach.  To compare
it with $D$, we must also retain curvature away from the base: a local
curvature value alone does not control the two-sided comparison problem.
The exponentially weighted supremum $J_0$ in the next proposition
records that curvature, with value $1$ assigned to the excluded set $X$.

\Needspace{5\baselineskip}
\begin{proposition}[Exact exponential comparison]
\label{prop:interface-exact-comparison}
There exist $\lambda>0$ and $C<\infty$, depending on the fixed
penalty and dimension but independent of $\delta\in(0,1/4]$, such that
\begin{equation}\label{eq:L5}
 D(z)\le C\bigl(J_0(z)+E(z)\bigr),\qquad
 E(z)\le C\bigl(J_0(z)+D(z)\bigr),\qquad z\in\Sigma\setminus X,
 \quad J_0(z)=\sup_{y\in\Sigma}e^{-\lambda|y-z|}\bar a(y)^2.
\end{equation}
\end{proposition}

We prove the proposition with explicit comparisons.  In particular,
no loss in the exponent $\sqrt2$ is absorbed into a constant.
The lower bound for $D$ comes from a positive kernel and the opposite
sheet.  For the upper bound, we construct an isolated solution on each
side of the same base graph.  Adding their equations cancels the terms
linear in the graph's mean curvature.  The remaining task is to compare
each artificial solution with the true phase while preserving the exact
exponential rate across a growing chamber.

\begin{lemma}[A half-space inverse]\label{lem:interface-halfspace-inverse}
Let $p_0>n$, $V_g=3g^2-1$, and
$L_g=-\Delta_y-\partial_t^2+V_g(t)$ on $\R^m\times(0,\infty)$,
with zero Dirichlet trace.  For some $b_0>0$, uniformly in
$0\le b\le b_0$, its inverse maps the norm
\[
 \sup_{(y,t)}e^{-b|y|}\|f\|_{L^{p_0}(B_1(y,t)\cap\{t>0\})}
\]
to the analogous uniformly local $W^{2,p_0}$ norm.  The inverse
is unique in this weighted class.
\end{lemma}

\begin{proof}
Let $v_0$ be the bounded positive half-line function supplied by
Lemma~\ref{lem:reach-half-line-barrier} with initial endpoint $0$.
It satisfies $(-\partial_t^2+V_g)v_0=1$,
$0<c\le v_0\le C$, and $|v_0'|+|v_0''|\le C$.
Its trace at $0$ is positive; the functions to which we apply the
inverse, rather than this comparison function, have zero trace.
Ground-state substitution, first for compact zero-trace functions,
gives
\[
 \int\bigl(|\nabla f|^2+V_g f^2\bigr)
 =\int v_0^2|\nabla(f/v_0)|^2+\int v_0^{-1}f^2
 \ge C^{-1}\|f\|_2^2.
\]
The energy inverse therefore exists for compactly supported sources.
For a source in a unit block, energy estimates followed by interior
and flat-boundary elliptic estimates and Sobolev bootstrapping bound
its $W^{2,p_0}$ norm locally by the source's $L^{p_0}$ norm.  These
bounds are uniform in the block location; only a fixed number of
bootstraps is required because $p_0>n$ is fixed.

Off a fixed enlargement of that block, compare its absolute value
with $Cv_0(t)e^{-c_1 r}$, where $r$ is distance to the block center.
For sufficiently small fixed $c_1>0$, this is a supersolution:
$L_gv_0=1$ dominates the terms bounded by
$C(c_1+c_1^2+c_1/r)v_0$.  Comparison is valid after division by
$v_0$, since the divided operator has positive zeroth coefficient
$v_0^{-1}$.  On unbounded domains it follows by exhaustion using
outer growing exponential barriers; alternatively, the energy
solution's vanishing exterior $L^2$ norm and local estimates give
the boundary condition at infinity.  Local elliptic estimates then
give exponential decay of the $W^{2,p_0}$ norms as well.
Partition an arbitrary source into disjoint unit blocks and sum
these estimates.  The sum converges locally absolutely because
$e^{b|y'|}\le e^{b|y|}e^{b|y-y'|}$ and $b<c_1/2$.
It provides the stated inverse and norm bound.  Here is an explicit
uniqueness argument for the full weighted class, including directions
with $t\to\infty$.  Let $R_*(x)=\sqrt{1+|x|^2}$, and write
$C_0=\sup v_0$, $C_1=\sup|v_0'|$.  Since
$|\nabla R_*|\le1$ and $\Delta R_*\le n$,
\[
 L_g(v_0e^{cR_*})
 \ge e^{cR_*}\bigl[1-2cC_1-ncC_0-c^2C_0\bigr].
\]
Choose $c_*>0$ so the bracket is at least $1/2$, and decrease $b_0$
to be at most $\min(c_1/4,c_*/2)$.  A homogeneous solution in the
weighted class satisfies $|u(y,t)|\le C_u e^{b|y|}$ by local Sobolev
embedding.  On the hemisphere $|x|=R$, the barrier
$\eta v_0e^{c_*R_*}$ dominates $|u|$ for every sufficiently large
$R$, since $c_*>b$.  On its flat boundary $u=0$.  The divided
maximum principle on half-balls, for both signs of $u$, therefore
gives $|u|\le\eta v_0e^{c_*R_*}$ everywhere.  Let $\eta\downarrow0$.
\end{proof}

\begin{proof}[Proof of Proposition~\ref{prop:interface-exact-comparison}]
Fix $z$ and write $B=J_0(z)+E(z)$.  We need only treat $B\to0$.
For the first inequality other values are trivial.  For the second,
if $E\le J_0$ there is nothing to prove; if $B$ is bounded below and
$E>J_0$, then $d$ is uniformly bounded and $D$ is uniformly bounded
below.  The latter assertion follows from single-profile compactness:
$D(z_j)\to0$ would make $U$ converge to a linear coordinate on every
fixed ball, forcing $d(z_j)\to\infty$ by the definition of the penalty.
This is also Lemma~\ref{lem:reach-small-defect}.

\emph{Step 1: ordered sheets on the comparison scale.}
Translate and rotate so $z=0$, $\nu(z)=e_n$.  Put
$P_0=A_0\log(1/B)$, where $A_0$ will be large and fixed.
After $A_0$ is chosen, take $\lambda$ so small that
$10\lambda A_0<1/2$.  Every zero in $B_{10P_0}$ is then good, and
\[
 |A|^2\le J_0(0)e^{10\lambda P_0}\le B^{1/2}.
\]
Indeed a bad zero there would give $J_0(0)\ge B^{10\lambda A_0}>B$.
At every such good zero $|\nabla v|\ge\sqrt{3/8}$, so uniform
elliptic estimates give uniform graph smoothness on a fixed scale.

The base sheet extends over a disk of radius at least $2P_0$ in its
tangent plane.  To see this directly, lift radial segments by tangent
projection.  Along a lift of length $O(P_0)$ the normal changes by
at most $O(P_0B^{1/4})=o(1)$, so the projection remains invertible
and the lift length remains controlled.  All lifts stay inside the
regular larger ball, and any terminal limit is regular.  Continuation
and local uniqueness give a smooth graph; the same argument works
for sheets based in a fixed small fraction of $B_{P_0}$.  Their
heights differ from their tangent planes by $O(P_0^2B^{1/4})=o(1)$.

Here is the ordering argument with the length scales recorded.
Fix a small angular parameter $\theta>0$.  A sheet based within
$\theta P_0$ has unoriented normal within $C\theta+o(1)$ of $e_n$.
Otherwise its almost affine disk of radius comparable to $P_0$
crosses the base disk transversally: along a line in the direction
of its tilt, its signed height relative to the base changes sign,
with magnitude bounded below by a constant times the tilt times
$P_0$, while both graph errors are $o(1)$.  The intermediate value
theorem would give an intersection, impossible for two different
regular zero sheets.  Thus all pieces in a smaller box are ordered
graphs over one horizontal disk.  There are finitely many in each
closed inner box: infinitely many would accumulate at a regular zero,
where the fixed local graph property permits only one piece.
Descriptions on overlapping disks agree or are disjoint.

Let the first graph above the base have positive-axis intercept
less than $\theta P_0/10$, if such a graph exists, and let $\ell$
be its shortest distance to $0$.  Its foot normal $n'$ satisfies
\begin{equation}\label{eq:interface-gap-angle}
 |n'-e_n|\le C\frac{\ell+o(1)}{P_0},\qquad
 \{n'\cdot x=\ell\}\text{ is its tangent plane at the foot}.
\end{equation}
For this stronger bound, apply the preceding nonintersection argument
out to the full horizontal scale comparable to $P_0$, now using
its height $\ell$ instead of $\theta P_0$.  The minimizing foot
is interior because the horizontal chart radius greatly exceeds
its intercept.  Choosing $\theta$ small puts the chord in the
constant part of $\Gamma$, and hence $d(0)\le G_1\ell$.

Choose box fractions $c_1=c'_0\theta$, $c_2=c_1/C_0$, with $c'_0$
sufficiently small and $C_0$ sufficiently large dimensional constants.
These two constants can be fixed before $\theta$ is decreased further.
In $|y|_\infty\le10c_1P_0$ the adjacent height differs from $\ell$
by at most a small fixed fraction of $\ell$, plus $o(1)$, by
\eqref{eq:interface-gap-angle}.  This box contains the whole adjacent
chamber and ample lateral buffer.  If no small intercept exists,
the analogous box contains no zero above the base up to height
$10c_1P_0$.  An intervening zero would extend through the axis to
a prohibited smaller intercept.  The same statements hold below.
All these chamber heights tend to infinity, since
$d(0)\ge\kappa^{-1}\log(1/B)$.

If $E(0)\ge J_0(0)$, then
$d(0)=\kappa^{-1}\log(1/E(0))=O(\log(1/B))$.
For $A_0$ sufficiently large it is much smaller than $\theta P_0$.
Every realizing contact lies in the preceding graph system.  A contact
on the base sheet, including a diagonal one, has value at least
$cB^{-1/4}\gg P_0$, by its curvature bound and the graph Taylor
formula.  Every other competing sheet has its shortest chord in the
plateau, while every chord to it has penalty at least $G_1$.
Consequently
\begin{equation}\label{eq:interface-real-gap}
 d(0)=G_1\ell_{\min},
\end{equation}
where $\ell_{\min}$ is the shortest distance to an adjacent graph.
Only adjacent graphs can be shortest, since a segment to a farther
graph crosses an intervening zero.

\emph{Step 2: the opposite sheet gives the lower interaction bound.}
For a fixed unit vector $n_0$, put $w=q(1-p\cdot n_0)\ge0$.
The identities for $q$ and $\nabla v$, together with
$q p=\sqrt2\nabla v$, give
\begin{equation}\label{eq:L6}
 (-\Delta+2)w=3qw-q^2D\ge qw\ge0.
\end{equation}
The last inequality uses
$D\le2(1-p\cdot n_0)$.
The positive kernel of $-\Delta+2$ is
\[
 \mathcal K(x)=\int_0^\infty(4\pi t)^{-n/2}
                       e^{-2t-|x|^2/(4t)}\,\dd t.
\]
For $r\ge1$ it is bounded above and below by positive constants
times $r^{-m/2}e^{-\sqrt2 r}$; the bounds follow by splitting this
integral about its minimum at $t=r/(2\sqrt2)$ and scaling the
central interval by $\sqrt r$.  Convolution represents bounded
solutions of this equation, by the maximum principle for their
bounded homogeneous difference.  In particular a nonnegative
subset of its source gives a lower bound.

Take $n_0=e_n$ and assume $E\ge J_0$.  On the opposite adjacent
sheet realizing \eqref{eq:interface-real-gap}, the oriented normal
is almost $-e_n$, because the signs of a continuous function alternate
between its regular zero graphs.  On a tangent disk of radius
$c\sqrt{\ell_{\min}}$ and a fixed thin neighborhood of this disk,
$q w$ has a fixed positive lower bound: $D<1/4$ at its zeros and
uniform local derivative bounds apply.  The distances to $0$ are
$\ell_{\min}+O(1)$, and its volume is comparable to
$\ell_{\min}^{m/2}$.  Thus \eqref{eq:L6} gives
\[
 D(0)\ge w(0)\ge c e^{-\sqrt2\ell_{\min}}=cE(0).
\]
This proves the second inequality of \eqref{eq:L5}, including its
exact exponential rate.  The remaining inequality requires an upper
bound for the deficit of the base slope; a one-sided curvature error
would be too large for that purpose.

\emph{Step 3: cancellation for two artificial half-space solutions.}
Write the base graph as $x_n=h(y)$, with $h(0)=0$ and
$\nabla h(0)=0$, and cut it off smoothly outside the preceding
comparison region at scale $P_0$.  Its extension can be chosen with
\[
 \|h\|_{C^2}=o(1),\qquad
 |\nabla h(y)|+|\nabla^2h(y)|
          \le C\sqrt{J_0(0)}e^{2\lambda|y|}.
\]
For an explicit bound independent of $P_0$, take the cutoff equal
to $1$ on $|y|\le P_0$ and supported on $|y|\le2P_0$.
The uncut graph has $|\nabla h|\le1$, so
$|(y,h(y))|\le2|y|$ and
$|\nabla^2h(y)|\le C\sqrt{J_0(0)}e^{\lambda|y|}$.
Radial integration from $h(0)=\nabla h(0)=0$ gives
\[
 |\nabla h(y)|\le C\lambda^{-1}\sqrt{J_0(0)}e^{\lambda|y|},
 \qquad
 |h(y)|\le C\lambda^{-2}\sqrt{J_0(0)}e^{\lambda|y|}.
\]
Derivatives of the cutoff contribute factors $P_0^{-1}$ and
$P_0^{-2}$, which are at most $1$ for $P_0\ge1$.  Thus the extended
graph obeys the displayed weighted estimate with a constant at most
$C(1+\lambda^{-1}+\lambda^{-2})$, fixed after $\lambda$, and no
factor growing with $\log(1/B)$.  Its unweighted $C^2$ smallness
follows separately from $P_0^2B^{1/4}\to0$.
Use this same extension $h$ on both sides.  In the coordinates
$x_n=h(y)\pm t$, $t>0$, seek positive solutions
$v^*_\pm=g(t)+\varphi_\pm$ with zero boundary.  Their equations are
\begin{equation}\label{eq:L7}
 \begin{split}
 L_g\varphi_\pm={}&|\nabla h|^2\partial_{tt}(g+\varphi_\pm)
 \mp2\nabla h\cdot\nabla_y\partial_t\varphi_\pm
 \mp\Delta h\,\partial_t(g+\varphi_\pm)\\
 &-3g\varphi_\pm^2-\varphi_\pm^3.
 \end{split}
\end{equation}
Lemma~\ref{lem:interface-halfspace-inverse} and the embedding
$W^{2,p_0}_{\rm uloc}\subset C^{1,\eta}$ make the right-hand side
followed by the inverse a contraction in a small unweighted ball.
The highest derivative terms have coefficients tending uniformly
to zero; the nonlinear terms have Lipschitz constants tending to
zero on that ball.  The source terms from $g$ tend to zero there.
This proves existence and uniqueness of small zero-trace solutions;
local elliptic bootstrapping gives smoothness.  Weighted estimates
in the same equation, with the small-coefficient terms absorbed,
then give
\[
 \|\varphi_\pm\|_{W^{2,p_0}(B_1(y,t)\cap\{t>0\})}
       \le C\sqrt{J_0(0)}e^{3\lambda|y|}.
\]
They are positive by $C^1$ smallness near $t=0$ and $C^0$ smallness
away from it.  They are at most $1$: a supremizing sequence with
value above $1$, translated in the interior and using uniform
regularity, gives a positive interior maximum incompatible with
$\Delta v^*=v^{*3}-v^*$.

Put $S=\varphi_++\varphi_-$ and $T_*=\varphi_+-\varphi_-$.
Adding \eqref{eq:L7} gives exactly
\[
 \begin{split}
 L_gS={}&|\nabla h|^2(2g''+S_{tt})
       -2\nabla h\cdot\nabla_y(T_*)_t-\Delta h\,(T_*)_t\\
       &-3g(\varphi_+^2+\varphi_-^2)
                       -(\varphi_+^3+\varphi_-^3).
 \end{split}
\]
The terms $\mp\Delta h\,g'$ cancel.  Every remaining source is
$O(J_0(0))$ in the weighted local $L^{p_0}$ norm with exponent
$10\lambda$: the coefficient factors have weight $2\lambda$,
the perturbations have weight $3\lambda$, and the cubic terms
also have a uniform small bound on one factor.  Taking
$10\lambda<b_0$ and using the inverse and boundary Sobolev embedding,
we conclude
\begin{equation}\label{eq:L8}
 |\partial_t(\varphi_++\varphi_-)(0,0)|\le C J_0(0).
\end{equation}

Thus the sum of the two artificial boundary slopes differs from the
planar value by a curvature-square envelope.  We next estimate the
change in each slope caused by the actual adjacent sheet, if present.

\emph{Step 4: comparing an artificial solution with the true chamber.}
We prove, on the upper side with a small adjacent intercept,
\begin{equation}\label{eq:L9}
 |\partial_n(v-v^*_+)(0)|
       \le C\bigl(e^{-\sqrt2\ell}+e^{-cP_0}\bigr).
\end{equation}
If no small intercept is present, the right side is $Ce^{-cP_0}$.
The lower-side assertion compares $-v$ with $v^*_-$.
All constants $c>0$ below are fixed before choosing $A_0$.

\emph{Middle-chamber deficit.}
First we establish a sharp bound at the middle of a chamber.  Set
$t_1=x_n$ and $t_2=\ell-n'\cdot x$.  In the central lateral box,
near $t_1=t_2$, both $1-v$ and $1-v^*_+$ satisfy
\begin{equation}\label{eq:L10}
 1-v,\ 1-v^*_+
       \le C\bigl(e^{-\sqrt2 t_1}+e^{-\sqrt2 t_2}+e^{-cP_0}\bigr).
\end{equation}
Here and below $v$ is positive in the chamber.
For the proof, subtract the uniform $o(1)$ graph errors from $t_1$
and $t_2$, and add the lateral face distances of a wider box.
Call all these inward unit affine distances $l_1,\ldots,l_k$.
Their positive intersection is contained in the phase, by ordering.
For either deficit $u=1-v$ or $u=1-v^*_+$,
\[
 (-\Delta+2)u=3u^2-u^3.
\]
Sign-ball clearing, as proved with \eqref{eq:phase-decay}, makes
$u$ uniformly as small as desired where $\min_jl_j\ge M$, for
large fixed $M$.  Put $a_0=\sqrt2$ and choose positive constants
$C_j$ with $C_je^{-a_0M}=\eta_0$, where $\eta_0>0$ is fixed small.
On this inner polytope the function
\[
 \mathcal W=\sum_{j=1}^k
           \left(C_je^{-a_0l_j}-kC_j^2e^{-2a_0l_j}\right)
\]
is positive and at most $k\eta_0$.  With $r_j=C_je^{-a_0l_j}$,
\[
 (-\Delta+2)\mathcal W=6k\sum_jr_j^2
       \ge6\Bigl(\sum_jr_j\Bigr)^2\ge3\mathcal W^2.
\]
Choose $\eta_0$ so $k\eta_0<1/8$, then $M$ so clearing gives
$u\le\eta_0/2$ on the boundary.  On a boundary face
$\mathcal W\ge\eta_0-k\eta_0^2\ge\eta_0/2$.
The difference equation has positive zeroth coefficient while
$u,\mathcal W$ are small, so the maximum principle gives
$u\le\mathcal W$.  The central point is a distance comparable to
$P_0$ from the lateral faces; the $o(1)$ shifts only change constants.
This proves \eqref{eq:L10}.

\emph{Potential comparison on the lower half of the chamber.}
Now use the domain between $x_n=h(y)$ and the affine midpoint
$t_1=t_2$, with $|y|_\infty<c_2P_0$.  In the case without a small
intercept, use instead its horizontal top at height $c_2P_0$.
The difference $u_*=v-v^*_+$ solves
\[
 (-\Delta+V_*)u_*=0,
 \qquad V_*=v^2+vv^*_++(v^*_+)^2-1.
\]
For some fixed $\gamma_1>0$,
\begin{equation}\label{eq:interface-potential-comparison}
 V_*(x)\ge V_g(t_1)-\eta(B)e^{-\gamma_1t_1},\qquad
 \eta(B)\longrightarrow0.
\end{equation}
We check both the fixed-depth and increasing-depth content of this
statement.  At each base point in the lateral box there is a tangent
ball in the positive phase whose radius tends uniformly to infinity.
It may be chosen much smaller than $c_1P_0$, the adjacent height,
and $B^{-1/4}$: the base Hessian bound keeps the ball above the base,
and the chamber and lateral buffers exclude any other zero from it.
For a concrete uniform choice, put $H_B=\ell$ when the small adjacent
intercept is present and $H_B=c_1P_0$ otherwise, and take
\[
 r_B=\bigl[\min\{c_1P_0,H_B,B^{-1/4}\}\bigr]^{1/2}.
\]
All three quantities in the minimum tend to infinity, so $r_B$ is
smaller than each by a factor tending to zero.  The adjacent height
over the larger box is bounded below by a fixed positive fraction of
$H_B$, and every base in the smaller box has a lateral margin
comparable to $c_1P_0$.  Finally $r_B B^{1/4}\to0$, so the Hessian
bound excludes a second contact of the tangent ball with the base.
The radius therefore works uniformly, including at bases on a lateral
face of the smaller box.
Lemma~\ref{lem:reach-phase-ball}, or its explicit inequality
\eqref{eq:C6}, gives $v\to g(t_1)$ on every fixed depth interval.
Indeed, along the inward normal at a base point, that inequality and
$|\nabla U|\le1$ give
$s-C(1+s)/r\le U\le s$ for fixed $s\ge0$, where $r$ is the radius
of the tangent phase ball.  These radii tend uniformly to infinity.
To pass from normal depth to the vertical coordinate $t_1$, use the
uniformly vanishing slope and height of the base graph and the same
Lipschitz bound for $U$.  Thus the convergence is uniform over all bases
in the buffered comparison box; no two-sided estimate for $H$ is used.
The artificial solution has the same convergence by construction,
and $h=o(1)$ on this box.  At increasing depths up to the midpoint
there is an explicit phase-ball radius independent of $\theta$.
Let $\eta_B=o(1)$ bound the signed affine errors of both sheets in
the larger box $|y|_\infty\le10c_1P_0$.  For a point in the closure
of the smaller lateral comparison box, put $s=t_1\ge1$.  We have
$t_2\ge s$ and $s\le\theta P_0/4$; the latter follows from
$\ell\le\theta P_0/10$ and \eqref{eq:interface-gap-angle}.
Set $c_{\rm ball}=\min(1/8,c'_0)$.
On $B_{c_{\rm ball}s}(x)$ both unit affine distances are at least
$7s/8>\eta_B$ for small $B$.  Its horizontal displacement is at
most $c'_0\theta P_0/4=c_1P_0/4$, so it remains inside the larger
box even when its center is on a lateral face of the smaller box.
It therefore lies in the true positive chamber.  It also lies above
the artificial graph.  Without a small intercept, $s\le c_2P_0$,
and the same ball stays below the empty-chamber height $10c_1P_0$.
Sign-ball clearing consequently gives
$1-v+1-v^*_+\le Ce^{-c_3t_1}$ with $c_3>0$ depending only on the
fixed $c_{\rm ball}$ and dimension, before $\theta$ is chosen.
Take $0<\gamma_1<\min(c_3,\sqrt2)$.  Splitting at a fixed large depth and then
using fixed-depth convergence proves \eqref{eq:interface-potential-comparison}.
This reasoning includes the domain with the artificial upper face.

\emph{The barrier at rate $\sqrt2$.}
We have bounded the error in the potential by a decaying function
with coefficient $\eta(B)\to0$.  We now absorb that error into
the comparison operator.  This is the step that preserves the exact rate.
Choose
a fixed small $\theta_0>0$; all $t_1$ in the domain are at least
$-\theta_0$ for small $B$.  Put $a_0=\sqrt2$ and
\[
 Y(t)=g'(t)\left[1+\int_{-\theta_0}^t g'(s)^{-2}\,\dd s\right],
 \qquad
 k'(t)=-Y(t)^{-2}\int_{-\theta_0}^tY(s)^2e^{-\gamma_1s}\,\dd s.
\]
Then $Y>0$, $(-\partial_t^2+V_g)Y=0$, and
$Y(t)\asymp e^{a_0t}$ as $t\to\infty$.
Consequently $|k'(t)|\le Ce^{-\gamma_1t}$ and its integral is
finite.  Choose $k(-\theta_0)$ so $1\le k\le C$, and set
$\phi=e^{-a_0t}Yk$.  Direct differentiation gives
\begin{equation}\label{eq:interface-shifted-barrier}
 \begin{gathered}
 0<c\le\phi\le C,\qquad |\phi'|\le Ce^{-\gamma_1t},\\
 -\phi''-2a_0\phi'+(V_g-2)\phi
       =e^{-a_0t}Y(t)e^{-\gamma_1t}
       \ge c e^{-\gamma_1t}.
 \end{gathered}
\end{equation}
These bounds also hold on the fixed compact initial interval.
If $c'=n'\cdot e_n$, then
\[
 (-\Delta+V_*)\bigl[e^{-a_0t_2}\phi(t_1)\bigr]
 =e^{-a_0t_2}
   \bigl[-\phi''-2a_0c'\phi'+(V_*-2)\phi\bigr].
\]
By \eqref{eq:interface-shifted-barrier} the bracket is nonnegative
when $|1-c'|$ is a sufficiently small fixed number and $B$ is
sufficiently small.  Its errors are bounded by
$C(|1-c'|+\eta(B))e^{-\gamma_1t_1}$.
This is the reason the exact exponential $e^{-\sqrt2t_2}$ survives.
The small geometric and potential errors have been absorbed in the
operator inequality before comparison.  They are not additive errors
in the interaction estimate and are never divided by its exponentially
small scale.

There are also coarse barriers for the remaining faces.
Take $v_0$ from Lemma~\ref{lem:reach-half-line-barrier} with
initial endpoint $-\theta_0$.  It is bounded above and below by
positive constants and satisfies $(-\partial_t^2+V_g)v_0=1$.
For any inward affine unit face distance $l$, the function
$v_0(t_1)e^{-c_4l}$ is a supersolution of $-\Delta+V_*$ when
$c_4>0$ is sufficiently small and $B$ is sufficiently small.
Indeed the new terms are bounded by $C(c_4+c_4^2+\eta(B))$,
whereas the undisturbed positive term is $1$.

The difference is zero at the bottom.  Multiples of these coarse
barriers dominate its absolute value at the lateral faces.  At the
midpoint, \eqref{eq:L10} and $t_1=t_2$ bound it by a multiple of
$e^{-a_0t_2}\phi(t_1)$ plus $e^{-cP_0}v_0(t_1)$.
The latter is itself a supersolution after multiplication by the
constant $e^{-cP_0}$.  In the case with an artificial top face use
its coarse barrier instead.  Comparison is legitimate despite the
possible sign of $V_*$: division by $v_0$ gives a strictly positive
zeroth coefficient, by \eqref{eq:interface-potential-comparison}.
The resulting bound on a fixed neighborhood of $0$ is
$C(e^{-a_0\ell}+e^{-cP_0})$, or $Ce^{-cP_0}$ without an adjacent
small intercept.  Boundary gradient estimates for the homogeneous
equation on the uniformly smooth graph, with zero boundary value,
prove \eqref{eq:L9}.

\emph{Step 5: choosing the constants and recovering the base defect.}
The artificial cancellation and the true-chamber comparison are now
proved separately.  To combine them uniformly, use the following order
for the constants.  First fix
the decay and comparison constants for $g$, the half-space inverse,
and the elementary phase-ball clearing.  The box construction above
has uniformly controlled aspect ratios when both fractions are
proportional to $\theta$, so choose $\gamma_1$ and then $\phi,v_0$
using those uniform constants.  Next take $\theta$ small enough for
the plateau condition and the drift error in
\eqref{eq:interface-shifted-barrier}; choose $c_1,c_2$ as specified.
The resulting coarse exponent $c>0$ is fixed.  Now choose $A_0$
large enough for the realizing-contact argument and $cA_0>1$.
Then choose $\lambda$ small enough for
$10\lambda A_0<1/2$ and $10\lambda<b_0$.
Finally decrease the upper threshold for $B$ to meet all small-error
conditions.  This order is independent of $\delta$.
In particular $e^{-cP_0}\le B$ with uniform constants.

At $0$, the two unsigned true derivatives in the inward height
coordinates are both $\partial_nv(0)$.  Equations \eqref{eq:L8}
and \eqref{eq:L9} show
\[
 |2\partial_nv(0)-2g'(0)|
       \le C\bigl(J_0(0)+E(0)+e^{-cP_0}\bigr)\le CB,
\]
because every adjacent gap used satisfies $d(0)\le G_1\ell$.
Since $D(0)=1-2\partial_nv(0)^2$ and $g'(0)=1/\sqrt2$,
this proves the first inequality of \eqref{eq:L5} and completes
the proposition.
\end{proof}

\subsection{An integral comparison}

The pointwise comparison contains the supremum $J_0$.  For the growth
argument we need an integral estimate involving $a^4$, $D^2$ and $E^2$
instead.  Uniform local area and derivative bounds let us replace the
supremum by a positive convolution, with a small $D^2$ error that can
be absorbed.  The constants must remain independent of the eventual
choice of the defect threshold and of the center of the spatial weight.

\begin{lemma}[Uniform area on good pieces]\label{lem:interface-good-area}
For $0<\delta\le1/4$, the area of $\Sigma\setminus X$ in any unit
ball is bounded by a dimensional constant, independently of $\delta$.
Consequently its area in an ambient ball of radius $r\ge1$ is at
most $C(1+r)^n$.  At a good zero whose distance from $X$ exceeds $2$,
there is a fixed-size graph neighborhood on which $D$ is comparable
to its value at that zero and $|\nabla_\Sigma A|\le C\sqrt D$.
\end{lemma}

\begin{proof}
At a good zero, $|\nabla v|\ge\sqrt{3/8}$.
Uniform second derivative bounds for $v$ give a fixed ball on which
its derivative along the initial normal is bounded below.  All zeros
in a smaller ball therefore form one graph, with uniformly bounded
slope and derivatives.  A separated selection of such smaller balls
has bounded cardinality inside a unit ball, by ordinary ambient
packing, and covers its good zeros.  Summing their graph areas gives
the claim.  Larger balls are covered by $C(1+r)^n$ unit balls.
The last assertions are local Harnack and the differentiated bounds
in Lemma~\ref{lem:diffuse-local}, applied on a fixed regular
graph contained in the good set.
\end{proof}

\Needspace{3\baselineskip}
\begin{lemma}[Weighted integral comparison]\label{lem:interface-integral}
Let $\mathcal X$ be a maximal unit-separated subset of $X$, with
$\mathcal X=\varnothing$ when $X=\varnothing$.  For an arbitrary
center $x_c$, set
$\chi_R(x)=\exp(-\sqrt{1+|x-x_c|^2}/R)$.
There are $C_D,R_D<\infty$, independent of
$0<\delta\le1/4$ and $x_c$, such that for $R\ge R_D$,
\begin{equation}\label{eq:L11}
 \int_{\Sigma\setminus X}(D^2+E^2)\chi_R^2\,\dd S
 \le C_D\left\{
 \int_{\Sigma\setminus X}(a^4+DE)\chi_R^2\,\dd S
             +\sum_{x\in\mathcal X}\chi_R(x)^2\right\}.
\end{equation}
\end{lemma}

\begin{proof}
Multiplying the first inequality of \eqref{eq:L5} by $D$ and the
second by $E$, and absorbing $DJ_0,EJ_0$ by Young's inequality,
gives $D^2+E^2\le C(J_0^2+DE)$.
We estimate $J_0^2$ without replacing a supremum by an unjustified
point value integral.  Let $w$ be good and $\dist(w,X)>2$.
On its fixed regular graph, for $w'$ within intrinsic distance $\mu$,
Lemma~\ref{lem:interface-good-area} gives
\[
 a(w)^4\le C a(w')^4+C\mu^4D(w)^2.
\]
Average on that small graph disk, and use Harnack and a fixed-size
graph disk to bound $D(w)^2$ by the integral of $D^2$ in $B_1(w)$.
For every sufficiently small fixed $\mu>0$ this yields
\[
 a(w)^4\le C_\mu\int_{B_1(w)\cap(\Sigma\setminus X)}a^4\,\dd S
       +C\mu^4\int_{B_1(w)\cap(\Sigma\setminus X)}D^2\,\dd S.
\]
The coefficient in front of the last integral is independent of
$\mu$.  Values at $w$ within distance $2$ of $X$, including bad
points where $\bar a=1$, are bounded by a constant and can instead
be charged to some point of $\mathcal X$ within distance $3$.
Since bounded shifts change $e^{-2\lambda|z-w|}$ by only a fixed
factor, taking the supremum over $w$ gives
\[
 J_0(z)^2\le C\int e^{-2\lambda|z-w|}
   \left(C_\mu a^4\,\dd S+C\mu^4D^2\,\dd S
                                  +\dd\#_{\mathcal X}\right)(w).
\]
The surface measures here are restricted to $\Sigma\setminus X$.
By Lemma~\ref{lem:interface-good-area} and
$\chi_R(z)\le e^{|z-w|/R}\chi_R(w)$, for $R\ge2/\lambda$,
\[
 \int_{\Sigma\setminus X}
   e^{-2\lambda|z-w|}\chi_R(z)^2\,\dd S_z
       \le C_\lambda\chi_R(w)^2.
\]
The constant is independent of the good threshold, by summation over
unit annuli and their polynomial area bound.  Tonelli now bounds
the integral of $J_0^2\chi_R^2$.  Choose $\mu$ small enough to
absorb the resulting $D^2$ term, after the constant in
\eqref{eq:L5} and $\lambda$ have been fixed.  This proves
\eqref{eq:L11}.  All integrals are finite: $D,E,a$ are bounded on
good zeros, the area has polynomial growth, and the separated net
has polynomial ambient growth.
\end{proof}

\subsection{Sharp constants in logarithmic windows}

The comparison above fixes the exponential rate, with unspecified
multiplicative constants.  On windows whose defect decreases at a
polynomial rate, the sheets are sufficiently close to parallel planes
to compute the coefficients as well.  We first make the geometry and
the additive distance errors precise; these errors must tend to zero
before they are inserted into an exponential.

\begin{lemma}[Uniform geometry on clean windows]
\label{lem:interface-log-windows}
Fix $\omega>0$, $C_0<\infty$, and a fractional margin $0<\theta_0<1$.
Suppose that $z_0\in\Sigma$ and
\[
 D(w)\le N^{-\omega}\qquad
 (w\in\Sigma\cap B_{N^\omega}(z_0)),\qquad N\longrightarrow\infty.
\]
The following conclusions hold uniformly for bases
$z\in\Sigma\cap B_{(1-\theta_0)N^\omega}(z_0)$.
Local zero pieces within $C_0\log N$ of $z$ are finitely many ordered
smooth graphs over $T_z\Sigma$.  Their domains contain a disk of
radius $cN^{\omega/8}$, with a further comparable graph buffer.
They have curvature at most $CN^{-\omega/2}$ and slopes and differences
of unoriented normals at most
\[
 C\bigl(N^{-\omega/8}\log N+N^{-3\omega/8}\bigr).
\]
Distinct graphs have Euclidean separation at least $c_0\log N$
where they participate in these neighborhoods.

For every fixed $k>0$, these descriptions extend over lateral disks
of radius $(\log N)^k$ and through every fixed multiple of $\log N$
in normal depth.  They contain every zero piece meeting the
corresponding smaller cylinder.  Projection to each participating
local graph is unique and smooth there.  Differences between its
Euclidean distance and its vertical distance, and differences between
its height increments and those of its tangent plane, are additive
errors bounded by
\[
 C_kN^{-c}(1+\log N)^{C_k}=o(1),\qquad c=c(\omega)>0.
\]
The same bounds hold uniformly as the center moves in the stated
fractional interior.  In particular the lateral domains include
$(\log N)^2$ boxes and disks of radius $\ell^{2/3}$ whenever
$\ell=O(\log N)$.  The polynomial buffers and unique projections make
these descriptions consistent on overlaps.

If a distinct graph is within $C_0\log N$ of the base $z$, then
\[
 d(z)=G_1\ell_{\min}(z),
\]
where $\ell_{\min}(z)$ is the shortest Euclidean distance from $z$
to an adjacent graph.  All thresholds may depend on the fixed
parameters and on $k$, but not on the solution or the base.
\end{lemma}

\begin{proof}
We establish graph continuation on a polynomial scale before using it
on any logarithmic scale.  Put
\[
 \kappa_N=CN^{-\omega/2},\qquad r=N^{\omega/8}.
\]
At every zero under consideration, $|A|\le\kappa_N$ and
$|\nabla v|$ has a fixed positive lower bound.  The fractional margin
contains all balls of radius $100r$ used below once $N$ is large.
Starting at any such zero, lift a straight radial segment in its
tangent plane by tangent projection.  As long as the lift has length
at most $10r$, its normal changes by at most $10\kappa_Nr=o(1)$.
The differential of projection therefore remains invertible, and
the lifted length is at most twice the projected length.  In
particular every segment of projected length at most $r$ continues
to its endpoint: a terminating sequence stays in the larger ball,
has a regular zero as its limit, and local inversion continues it.
The lifting equation depends smoothly on the segment endpoint.
Its endpoints thus form one smooth graph over the full tangent disk;
local uniqueness identifies it with the original zero piece.
Starting with a larger fixed multiple of $r$ supplies the stated
buffer.  Taylor's formula on these graphs gives the affine error
\[
 C\kappa_Nr^2=O(N^{-\omega/4}).
\]

We next put nearby pieces over the same plane.  Two such graph disks
whose base points are at distance $b=o(r)$ have unoriented normal
angle at most
\begin{equation}\label{eq:interface-clean-angle}
 C\left(\frac b r+\kappa_Nr\right).
\end{equation}
Indeed, if the angle were larger, a line in one tangent plane in the
direction of maximal tilt would make its signed distance from the
other tangent plane take opposite signs, with magnitude greater
than both affine errors, inside a fixed smaller disk.  Points of
that graph project into the interior of the other disk.  The signed
graph height, measured in the second disk's normal direction, then
vanishes by the intermediate value theorem.  At this intersection
the two normals would still differ by more than their variations
$C\kappa_Nr$, contradicting the unique tangent plane of the regular
level set.  This argument also excludes differently tilted pieces
of a single connected component.  Taking $b=O(\log N)$ proves the
claimed common-plane slope bound.  Over a common smaller disk the
graphs either agree or are disjoint, and so are ordered.

For completeness, no new sheet can enter unnoticed through the side
of a logarithmic-power box.  A zero $w$ with lateral coordinate
$|y|\le(\log N)^k$ and height at most $C_0\log N$ first has
$b=O((1+\log N)^{\max(k,1)})$ in
\eqref{eq:interface-clean-angle}.  This already gives a slope tending
to zero polynomially and permits its graph to continue to the normal
axis through $z$.  Its height changes on that continuation by at most
\[
 C\left(\frac{(1+\log N)^{\max(k,1)}}r+\kappa_Nr\right)
       (1+\log N)^k=o(1).
\]
Its intercept is therefore within $(C_0+1)\log N$ of $z$ for large
$N$.  Applying \eqref{eq:interface-clean-angle} again at this
intercept gives the refined slope bound already stated.  Enlarging
the fixed depth constant in a buffered description therefore includes
all pieces in the smaller box.  There are only finitely many in each
closed inner cylinder: otherwise zeros on distinct graphs accumulate
at a regular zero, whose uniform local graph neighborhood contains
only one zero piece.  This is local finiteness, and imposes no global
bound on the number of sheets.

The same estimates give the required additive accuracy.  On a disk
of radius $L=(\log N)^k$, the error from a graph's own tangent plane
is at most $C\kappa_NL^2$.  If
$\vartheta_N=C(\log N/r+\kappa_Nr)$ denotes the common-plane slope
bound, changing from its tangent plane to the common horizontal
plane costs at most $C\vartheta_NL+C\kappa_NL^2$ in height.
Both quantities are bounded by a negative power of $N$ times a
fixed logarithmic power.  They tend to zero as numbers, rather than
only relative to the gap.

To check projections, write one graph as $x_n=f(y)$.  For a fixed
point $x$, half its squared distance to $(y,f(y))$ has Hessian
\[
 \Id+\nabla f\otimes\nabla f-(x_n-f)D^2f.
\]
On a buffered disk relevant to logarithmic depth this is positive
definite, with eigenvalues $1+o(1)$: the slope tends to zero and
$|x_n-f||D^2f|=o(1)$.  The same bound holds out to a fixed fraction
of the polynomial disk, since its height variation is at most
$C\vartheta_Nr+C\kappa_Nr^2=O(\log N)$.
A minimizing point cannot occur at that outer boundary, whose
horizontal distance is polynomially large, whereas the vertical
competitor has distance $O(\log N)$.  Strict convexity gives a unique
foot, and the implicit function theorem gives its smooth dependence
on $x$.  The stationary equation bounds the horizontal displacement
of the foot by $C\vartheta_N\log N$ and the difference between true
and vertical distances by $C\vartheta_N^2\log N$, after harmless
enlargement of the common slope bound.  These are the asserted
additive errors.  This projection is to a specified local graph, not to an entire
connected component.  Every foot that can compete for the nearest
zero lies within logarithmic distance, and hence belongs to the
buffered family by the preceding completeness argument.

To prove separation, set $r_0=c_0\log N$ with $c_0>0$ sufficiently
small.  Starting at any participating zero, the exponential bound in
\eqref{eq:diffuse-local} gives
$D\le CN^{-\omega/2}$ throughout $B_{4r_0}$.
Consequently $|T|\le CN^{-\omega/4}$ there.  Integration shows that
$p$ remains uniformly close to the initial unit normal, so its
component in that direction is positive.  All zeros in a smaller
cylinder then form a single graph.  In conjunction with the common
graph description, this excludes a distinct sheet at distance
less than a smaller fixed multiple of $r_0$.

Finally assume a distinct sheet lies within $C_0\log N$ of $z$.
Its nearest foot is interior to the polynomial graph disk, and the
shortest chord is normal to that graph.  The common normal estimate
puts its angle at $z$ in the plateau of $\Gamma$, for large $N$.
It supplies a penalized competitor of size $O(\log N)$.
Every better competitor has chord length $O(\log N)$ because
$\Gamma\ge1$, so lies in the same graph system.  A chord of length
$l$ to the base graph has normal displacement at most
$C\kappa_N l^2$.  Since $\Gamma(s)\ge1/s$, its penalized length
is at least $c\kappa_N^{-1}\gg\log N$, also in a diagonal limit.
For every other graph all chords cost at least $G_1$ times their
length, and its shortest chord attains this factor.  Only adjacent
graphs can minimize, since a chord to a farther one crosses an
intervening zero first.  This proves the exact identity for $d$.
All graph identifications and projections on overlapping inner
windows agree by their local uniqueness; all construction scales
fit within the original fractional margin.
\end{proof}

\begin{proposition}[Sharp interaction and middle-phase tail]
\label{prop:interface-sharp-gaps}
Under the hypotheses of Lemma~\ref{lem:interface-log-windows},
for each $\eps>0$ there is $P_\eps<\infty$ such that every fixed
$P\ge P_\eps$ has the following properties when $N$ is sufficiently
large.  If a different sheet is within $C_0\log N$ of the base, then
\begin{equation}\label{eq:L12}
 I_P(z)\ge(32-\eps)e^{-\sqrt2\ell_{\min}(z)}.
\end{equation}
If $x$ lies between two adjacent sheets, their distances $s_i,s_j$
from $x$ are both at most $C_0\log N$, they are the nearest sheets,
and $|s_i-s_j|\le C_1$ for a fixed $C_1$, then
\begin{equation}\label{eq:L13}
 q(x)\le(4+\eps)
       \left(e^{-\sqrt2s_i}+e^{-\sqrt2s_j}\right).
\end{equation}
Largeness of $N$ may depend on all the fixed parameters, including
$P,C_1,\eps$.  The window geometry and conclusions are uniform in
the solution and center.
\end{proposition}

\begin{proof}
For \eqref{eq:L12} let $\ell=\ell_{\min}$ and fix
$n_0=\nu(z)$ in \eqref{eq:L6}.  On the column $z+s\nu(z)$,
$|s|\le P$, \eqref{eq:interface-normal-accuracy} and local Harnack
imply
\[
 D=2(1-p\cdot n_0)-|p-n_0|^2,
 \qquad qD\ge(2-o(1))w.
\]
In fact $|p-n_0|^2=O_P(D(z)^2)$ whereas $D$ is comparable to
$D(z)$ there, so the last error is uniform and relative.
At the opposite foot use a tangent disk of radius $\ell^{2/3}$
and bounded depth $|t|\le Q$ in its affine normal coordinates.
The true orientations are opposed, and all profile and affine
errors on this patch tend uniformly to zero by
Lemma~\ref{lem:interface-log-windows} and
\eqref{eq:interface-normal-accuracy} at its bases.  Thus the exact
source on the right of \eqref{eq:L6} satisfies
\[
 3qw-q^2D\ge(6-o(1))G(t)^2,\qquad |t|\le Q.
\]
The constant is $6$, since $w\to2G$ there.
For fixed $P,Q$ the comparison is relative, as $G$ has a positive
minimum on the compact depth interval.

Integrating the heat-kernel formula for $\mathcal K$ in its
horizontal variables gives exactly
\[
 \int_{\R^m}\mathcal K(y,h)\,\dd y
       =\frac1{2\sqrt2}e^{-\sqrt2|h|}.
\]
Truncating this integral to $|y|\le\ell^{2/3}$ changes it by a
relative $o(1)$ for $h=\ell+O(P+Q)$: the excess distance outside
that disk is at least $c\ell^{1/3}$, and the kernel bounds used
above control the remaining polynomial volume factor.
Polynomial-small graph and angular errors contribute only $o(1)$
additive errors to distances on this disk, so do not change the
exponential factor by more than $1+o(1)$.
It follows, with signs depending on the side but irrelevant by
symmetry of $G$, that
\[
 \begin{split}
 \frac{I_P(z)}{e^{-\sqrt2\ell}}
 &\ge \frac{2-o(1)}\sigma\frac{6-o(1)}{2\sqrt2}
 \left(\int_{-P}^{P}e^{\sqrt2s}G(s)^2\,\dd s\right)
 \left(\int_{-Q}^{Q}e^{\sqrt2t}G(t)^2\,\dd t\right).
 \end{split}
\]
Choose $P,Q$ large and fixed, then $N$ large.  The elementary
substitution $r=e^{\sqrt2t}$ gives
\[
 \int_\R e^{\sqrt2t}G(t)^2\,\dd t=2\sigma.
\]
The limiting coefficient is therefore
\[
 \frac2\sigma\frac6{2\sqrt2}(2\sigma)^2=32,
\]
which proves \eqref{eq:L12}.

The complementary estimate concerns the value of the phase between
the same sheets.  Here the coefficient comes from the one-dimensional
tail $1-g(s)\sim2e^{-\sqrt2s}$ and from $q\le2(1-v)$ in a positive
phase.  The lateral boundary must be placed farther away than the
logarithmic sheet separation so that its contribution is negligible
relative to the principal tails.

For \eqref{eq:L13}, change the sign of $v$ if necessary so the
phase between the sheets is positive.  Separation and
$|s_i-s_j|\le C_1$ imply $\min(s_i,s_j)\to\infty$, so the point
$x$ belongs to the fixed-depth inner polytope used below once $N$
is sufficiently large.  Use the two tangent planes
at the perpendicular feet from $x$, together with lateral faces
at scale $(\log N)^2$.  Lemma~\ref{lem:interface-log-windows}
provides these graphs with polynomial buffers, no intervening
zeros, and affine errors $o(1)$.  Subtract these errors from the
two face distances, producing an inner polytope as in the proof
of \eqref{eq:L10}.  Its distances at $x$ are $s_i+o(1)$ and
$s_j+o(1)$.

In that proof the two principal coefficients $C_j$ can now be
chosen equal to $2+\eta$, with arbitrarily small fixed $\eta>0$.
Indeed at every point of their inner boundaries $l_j=M$, including
intersections with lateral faces, the true perpendicular projection
is unique and remains in the polynomially buffered clean window.
The signed distance to its sheet is $M+o(1)$: the affine error on
the entire $(\log N)^2$ box is polynomial-small, including its
gradient, and the squared-distance Hessian in
Lemma~\ref{lem:interface-log-windows} is positive definite.
Fixed-depth profile convergence at these feet is uniform, so the
solution converges there to $g(M+o(1))$.  Since
$1-g(M)=2e^{-\sqrt2M}+O(e^{-2\sqrt2M})$, choose $M$ large
so that
\[
 (2+\eta)e^{-\sqrt2M}
       -k(2+\eta)^2e^{-2\sqrt2M}>1-g(M)
\]
with fixed positive slack, and then $N$ large enough to retain it.
Keep the coarse fixed coefficients on the lateral faces.  Each
summand in the barrier is nonnegative; thus intersections of faces
cause no difficulty in the boundary comparison.  The same algebra
$(-\Delta+2)\mathcal W\ge3\mathcal W^2$ applies with these unequal
coefficients.  At $x$, the lateral contribution is
$O(e^{-c(\log N)^2})$, hence negligible relative to either of the
two exponentials, since $s_i,s_j\le C_0\log N$.
We obtain
\[
 1-v(x)\le(2+2\eta)
       \left(e^{-\sqrt2s_i}+e^{-\sqrt2s_j}\right).
\]
Finally $q=(1-v)(1+v)\le2(1-v)$; take $\eta$ sufficiently small
in terms of $\eps$ to prove \eqref{eq:L13}.
\end{proof}


The two sharp coefficients concern different quantities: $32$ bounds
normal defect mass from below, while $4$ bounds the phase tail from
above.  The radial test will compare them on the same adjacent gap.
The fourth-order estimate proved next first supplies the control of
rough regions needed to make that comparison globally.
